\chapter{信息几何动力学：统计状态、投影更新与迁移边界}
\label{chp:info_geometry_dynamics}

前两章给出了单元内部的状态算子以及多单元之间的耦合、冲突和级联机制，本章进一步回答一个数值问题：当智能系统用概率分布、协方差、相似核或可学习参数表示不确定状态时，更新方向应怎样依赖表示的几何，而不是依赖任意坐标尺度。我们把信息几何限定为统计模型上的计算工具，不把所有认知状态预设为黎曼流形，也不从测地线、对偶联络或自然梯度直接推出思维机制。讨论依次定义正定矩阵状态、统计流形的对偶结构、自然梯度的连续与离散实现，并重建原章关于模糊匹配、类比迁移和组合泛化的案例。最终结论是条件性的：几何可以规定哪些差异应被视为等价、怎样作最小改动以及误差如何传播，但泛化是否成立仍取决于任务分布、表示接口、因果结构和外部验收。

\section{正定矩阵状态：对象、度量与可计算更新}
\label{sec:section_3075}

我们先区分三类经常被混用的对象。局部活动向量记为$h\in\mathbb R^n$，它可以表示当前注意、特征或工作状态；随机变量$H$的协方差记为$P=\mathbb E[(H-\mu)(H-\mu)^\top]\in\mathbb S_+^n$；参数化概率族$p_\theta(y)$的Fisher信息记为$G(\theta)\in\mathbb S_+^d$。$P$度量状态样本的二阶离散程度，$G$度量参数微小变化造成的分布可区分度，二者作用空间和维数通常不同。只有在高斯位置族、线性观测且参数正好对应均值等专门条件下，协方差与Fisher信息的逆才可能形成简单关系。因此我们不把任意“思维协方差”宣布为Fisher矩阵之逆，也不把矩阵迹、熵、计算代价和硬件焦耳能耗混成一个量。

为了稳定计算，设$P\in SPD(n)=\{P=P^\top\mid v^\top Pv>0,\ \forall v\neq0\}$，其分量单位由$H$的分量单位乘积决定。若不同特征使用不同单位，必须先通过尺度矩阵$S$归一化为$\widetilde P=S^{-1}PS^{-\top}$，否则一个数值较大的坐标会主导距离。$SPD(n)$是开凸锥而不是一般意义上的矩阵李群；它可看作$GL(n)/O(n)$的商空间。其切空间$T_PSPD(n)$由对称矩阵组成。采用仿射不变度量时，我们定义
\begin{equation}
g_P(X,Y)=\operatorname{tr}(P^{-1}XP^{-1}Y),\qquad X,Y\in T_PSPD(n),
\label{eq:spd_metric_dynamics}
\end{equation}
相应距离为$d_{\mathrm{AI}}(P,Q)=\|\log(P^{-1/2}QP^{-1/2})\|_F$。该度量在合同变换$P\mapsto APA^\top$、$A\in GL(n)$下不变，含义是可逆线性重参数化不会改变距离；它并不表示所有信号增益都无影响，更不自动保持逻辑关系。若缩放会改变噪声模型、饱和区间或任务损失，那么缩放就是实质变化，不能被“不变性”抹去。

在这一定义下，$P_0$到$P_1$的一条最短测地线是
\begin{equation}
P(s)=P_0^{1/2}\left(P_0^{-1/2}P_1P_0^{-1/2}\right)^sP_0^{1/2},\qquad s\in[0,1].
\end{equation}
它解决的是两个已知正定矩阵之间的插值问题，不是“快思考”的普遍轨迹。实际推理还会改变目标、输入、结构和维数，未必沿固定流形的最短路。以在线状态估计为例，系统可用指数映射把切向更新$V_k$送回流形：$P_{k+1}=\operatorname{Exp}_{P_k}(\eta_kV_k)$；也可采用保证正定的Cholesky参数$P=LL^\top+\varepsilon I$。前者保持几何含义但矩阵分解昂贵，后者便于自动微分却依赖坐标选择。工程设计需报告$n$、条件数、正则量$\varepsilon$、更新步长和分解误差，而不是笼统要求所谓TPU--G执行“李群积分”。

正定矩阵还可作为局部结构网络与全局分布表示之间的接口。局部网络保留对象身份、角色边和可追溯证据；从一组带身份的激活样本可估计任务相关$P_q$，供全局不确定性、耦合方向或异常程度读出。反向控制只能把$P_q$生成的增益、门控或采样预算写回具有明确身份的节点，不能把相关矩阵当成对象关系图。两个特征高协方差可能来自共同输入，删去图中的一条边不一定改变它；一条关键因果边也可能因样本稀少而没有显著协方差。验证协议因此至少包括合同变换一致性、正定性、条件数、留出预测误差和结构干预响应。只有这些测试共同通过，矩阵状态才是有用的计算容器，而不是把“高维纠缠”换成另一种命名。

矩阵状态的时间更新还要区分样本变化与估计器变化。若在窗口$W_k$内用样本协方差估计$P_k$，窗口滑动会同时删除旧样本、加入新样本；$P_{k+1}-P_k$不只是被表示对象的动力学，也包含采样误差和窗口规则。采用指数遗忘可提高响应速度，却降低有效样本量；扩大窗口可改善条件数，却可能掩盖突变。系统应同时记录窗口长度、有效样本量、缺失值处理和置信区间，并用变化点检验决定是否仍可在同一流形上连续更新。若表示维度因对象新增而改变，应先通过显式嵌入或边缘化建立跨维接口，不能直接计算两个不同维矩阵的测地距离。这个限制尤其适用于开放世界Agent：对象数量、工具集合和任务变量会动态改变，固定$SPD(n)$只能描述一个模式区间内的统计状态，而不是完整生命周期的唯一空间。

\section{对偶联络与投影：预测、校准及其成立条件}
\label{sec:section_6511}

考虑严格正的概率分布族$\mathcal M=\{p_\theta(y)\mid\theta\in\Theta\subseteq\mathbb R^d\}$。若$p_\theta$二阶可微且可交换积分与求导，Fisher度量定义为$G_{ij}(\theta)=\mathbb E_\theta[\partial_i\log p_\theta(Y)\partial_j\log p_\theta(Y)]$。在指数族$p_\theta(y)=\exp(\theta^\top T(y)-\psi(\theta))b(y)$中，自然参数$\theta$和期望参数$\eta=\mathbb E_\theta[T(Y)]=\nabla\psi(\theta)$构成Legendre对偶坐标。指数联络在$\theta$坐标中平坦，混合联络在$\eta$坐标中平坦；“平坦”是指定统计族和联络下的性质，不表示现实推理没有曲率、冲突或模型错设。

对偶性通过凸势$\psi$诱导的Bregman散度表达。令$D_\psi(\theta_1,\theta_2)=\psi(\theta_1)-\psi(\theta_2)-\nabla\psi(\theta_2)^\top(\theta_1-\theta_2)$。当$R$是把$P$投影到适当的仿射子流形上、投影方向和散度方向与该子流形的平坦结构匹配时，才有广义毕达哥拉斯关系$D(P,Q)=D(P,R)+D(R,Q)$。任取三个分布并不满足这个等式，KL散度也不对称。原章把任意预测流和数据流称为黎曼正交，会把条件命题扩大为普遍定理。我们改为可检验的模型假设：在一次更新窗口内模型族固定、投影解存在且唯一、约束集在所用坐标中凸，并明确使用正向还是反向KL；条件不满足时只报告数值最优解，不能声称正交分解。

在预测—校准循环中，先验状态$q_k$经动力模型$T_k$得到预测分布$q^-_{k+1}=T_k\#q_k$，符号$\#$表示分布推前。观测$z_{k+1}$到达后，以似然$\ell(z_{k+1}\mid y)$和可行集$\mathcal C_{k+1}$构造校准问题
\begin{equation}
q_{k+1}=\arg\min_{q\in\mathcal C_{k+1}}
\left[D_{\mathrm{KL}}(q\|q^-_{k+1})-\beta_k\,\mathbb E_q\log \ell(z_{k+1}\mid Y)+\lambda_k\Omega(q)\right].
\label{eq:prediction_calibration_projection}
\end{equation}
第一项限制相对预测的改动，第二项吸收观测，第三项表达复杂度、安全或资源约束；$\beta_k$和$\lambda_k$的单位必须使各项可比较。若$\mathcal C$只由数据“定义”，却没有写明噪声、异常值和权限，投影会把错误观测强制变成事实。故现实校准要保留来源、时间、传感器状态和拒绝选项。目标下降说明该优化问题得到改善，不说明外部世界已经被正确识别。

这一循环可以作为TDCI式工程过程的几何实现，但不是所有智能系统的唯一机制。AgentOS可把$h(t)$中的候选分布视为$q$，把$E$中的带来源事件形成似然，把$\Theta$中的结构约束形成$\mathcal C$，由SPC提供任务相关压缩，由TCE调节$\beta_k$与探索强度，CCM在条件数或残差过大时触发降阶，Boundary限制观测和结构写入，Offloading将完整证据写入外部记忆$M_e$。粒子滤波、集合估计、逻辑真值维护或非概率控制器也能实现预测和校准。选择何种表示取决于多模态性、实时预算、可追溯要求与错误代价，不能从对偶联络的优雅性推导工程唯一性。

离散实现还要处理异步和模型切换。若目标、约束、数据度量或模型族随时间改变，第$k$步应记录版本$(T_k,\mathcal C_k,G_k,\ell_k)$；不能把旧切空间的增量直接应用到新坐标。对延迟观测，应回溯到对应历史状态重算或使用显式延迟模型，而不是在当前状态上假装正交投影。验收可用校准误差、对数评分、拒绝率、投影残差、回滚成功率以及外部任务结果。反例是模型族遗漏真实模式：即使每一步都精确完成KL投影，系统仍会稳定地收敛到错误近似。几何保证的是给定模型内部的关系，不负责替代模型批判。

投影方向还对应不同的错误偏好。最小化$D_{\mathrm{KL}}(q\|p)$通常倾向选择$p$支持集中的高密度模式，而最小化$D_{\mathrm{KL}}(p\|q)$更重视覆盖$p$给出的各个模式；在多峰召回或安全监测中，两者会给出实质不同的结果。模式寻优可能生成清晰但漏项的解释，质量覆盖可能保留低概率风险却增加计算负担。我们不能以“最小惊奇”代替方向选择，而应由任务损失说明漏报和误报的代价。若约束集合非凸，投影还可能有多个局部解，初始化与求解器就成为模型的一部分。此时应保留候选解、目标间隙和停止准则，并在外部证据到达后允许重新开启分支。所谓预测与修正的分离，是一种便于审计的接口设计；在端到端循环网络中两者可能同时发生，仍需用干预或消融确定各自贡献。

\section{自然梯度动力学：不变方向、离散条件与计算代价}
\label{sec:section_4509}

设待优化参数$\theta\in\Theta\subseteq\mathbb R^d$决定分布$p_\theta$，目标$L(\theta,A,u,t)$依赖参数、结构$A$、外部输入$u$和显式时间。普通梯度$\nabla_\theta L$依赖坐标的欧氏尺度；在Fisher度量$G(\theta)$正定时，自然梯度$\widetilde\nabla L=G(\theta)^{-1}\nabla L$由局部最速下降问题得到：在二阶KL近似$\frac12\delta\theta^\top G\delta\theta\leq\epsilon$下，最小化一阶目标变化$\nabla L^\top\delta\theta$，拉格朗日条件给出$\delta\theta\propto-G^{-1}\nabla L$。这是一个数学结果；“最速”相对于所选度量和局部近似，而不是对墙钟时间、训练算力或任务成功率的无条件最优。

连续模型写为$\dot\theta=-G(\theta,A,u,t)^{-1}\nabla_\theta L(\theta,A,u,t)$。沿实际轨迹的目标导数必须完整展开：
\begin{equation}
\frac{dL}{dt}=-\nabla_\theta L^\top G^{-1}\nabla_\theta L
+\left\langle\partial_A L,\dot A\right\rangle
+\left\langle\partial_u L,\dot u\right\rangle+\partial_tL.
\label{eq:natural_gradient_total_derivative}
\end{equation}
只有结构和输入固定、目标无显式漂移且$G$正定时，第一项才保证$L$不增。即使$L$下降，也只说明内部代理目标改善，不等于现实任务成功。若$G$奇异，意味着某些参数方向在当前分布下不可辨识；直接求逆会放大噪声。工程上可使用$G_\rho=G+\rho I$、截断特征值、块对角近似或共轭梯度求解，并报告阻尼$\rho$对不变性和偏差的影响。

离散更新为$\theta_{k+1}=\Pi_{\Theta_k}^{G_k}(\theta_k-\alpha_k\widehat G_k^{-1}\widehat g_k)$，其中$\widehat g_k$和$\widehat G_k$是批数据估计，$\Pi$是按约束度量定义的投影。若局部自然梯度向量场在相关区域Lipschitz且目标在该度量下局部光滑，足够小的$\alpha_k$可获得局部下降；随机估计还需控制方差与偏差。大步长会越出二阶KL近似范围，结构事件会改变维数，过时的预条件器会把更新送向错误方向。因此每步应检查实际KL半径、预测下降与实际下降之比、可行性和数值条件数，超限时缩步、重估度量或回滚。矩阵指数或retraction可用于流形参数，但必须区分精确指数映射与只保持一阶性质的廉价回缩。

“确定知识更难改变”不是由$G$大直接推出的普遍规律。Fisher信息大表示某参数方向的小变化容易改变预测分布；乘以$G^{-1}$会缩小该坐标更新，但目标梯度、先验、数据量和参数化同时决定最终变化。冗余参数可能具有很小特征值却不对应“模糊知识”，稀有但关键事件也可能在经验Fisher中几乎不可见。若要表达记忆刚性，应另定义保持代价、版本承诺、证据量或安全约束，并观察反事实更新后的性能，而不能把曲率单独解释为心理惯性。

计算代价是自然梯度能否使用的现实边界。完整$d\times d$矩阵需要$O(d^2)$存储，直接分解约为$O(d^3)$；大型模型通常只能使用Kronecker、对角、低秩或隐式乘积近似。这些近似改变了度量，必须以相同预算下的收敛速度、最终校准、鲁棒性和硬件能耗比较，而不是只看训练损失步数。对AgentOS，TCE可以把自然梯度视为一种控制候选，CCM根据预算和病态程度在全矩阵、块近似与普通梯度之间切换；$\Theta$的持久结构更新需比$h(t)$的瞬时状态更新采用更严格的证据和回滚门槛。这是工程分层，并非信息几何要求所有层共享同一参数流形。

自然梯度的坐标不变性也必须精确表述。若两套参数通过光滑双射重参数化并使用同一分布族的精确Fisher度量，连续自然梯度向量表示同一条分布轨迹；有限步长、阻尼、动量、裁剪和近似矩阵都会破坏这种严格等价。两个实现即使在第一步方向一致，累积舍入和批次噪声也可能使长期轨迹分离。工程验证可选择一个已知可逆的参数变换，比较两端每步分布KL、目标变化和约束违反量；只比较参数欧氏距离没有意义。若系统在结构编辑后更换参数族，应把编辑视为离散重置映射，先将可保留的预测和身份关系迁移，再重新估计度量。这样才能把连续优化的结论限定在有效区间内，避免由局部不变性推出整体同构或终身稳定。

\section{泛化与迁移：结构保持、吸引域和反例检验}
\label{sec:geometric_generalization_physics}

泛化不是几何动力学的必然，而是系统在指定训练信息之外仍满足任务准则的经验性质。设训练环境分布为$P_{\mathrm{tr}}$，目标环境集合为$\mathcal Q$，模型$f$的损失为$\ell(f(x),y)$，则我们关心$R_Q(f)=\mathbb E_Q\ell(f(X),Y)$以及最坏情形$\sup_{Q\in\mathcal Q}R_Q(f)$。几何结构的作用是限定$\mathcal Q$中哪些变化应保持预测、哪些变化需要重估；如果不说明环境变化集、任务不变量和允许误差，“处理未见样本”没有可证内容。原章的雨水入河、象棋迁移到战争、红与球组合等例子可以保留，但必须分别落到鲁棒吸引、关系映射和组合接口，而不能由最小作用量、规范协变或函子名称宣布成功。

第一类机制是局部鲁棒与吸引域。设离散状态动力学$x_{k+1}=F_q(x_k)$在任务$q$下存在吸引集合$\mathcal A_q$，若对初值集合$B_q$有$d(F_q^k(x),\mathcal A_q)\leq C\lambda^kd(x,\mathcal A_q)$、$0<\lambda<1$，则$B_q$内扰动会衰减。这可以解释“雨水落在邻近位置仍汇入同一河道”式模糊匹配，但只对已证明或估计的吸引域成立。落在分水岭另一侧、遭遇结构改变或受到对抗扰动时，近邻样本会走向不同结果。验证应测扰动半径、收缩率、类别边界和失效样本，而不是比较两条完整路径的KL并假定其趋零。几何距离小也不保证任务标签相同，任务标签相同也不要求内部路径同构。

第二类机制是跨域关系映射。令源域结构$G_s=(V_s,E_s,\tau_s)$和目标域结构$G_t$分别记录对象、关系与约束，迁移映射$\phi:V_s\to V_t$只有在一组任务相关关系上满足$\tau_t(\phi(a),\phi(b))\approx\phi_E(\tau_s(a,b))$，并且行动后果、边界条件和价值准则也经过校准时，才能复用策略。把“围魏救赵”从战争映射到商战，至少要对应资源约束、对手响应、时间窗和合法性；共享“趋利避害”不足以证明动力学等效。棋子可牺牲不意味着人员可牺牲，零和棋局也不等于开放市场。规范或坐标协变至多表达描述变化下某些关系保持，不能从相同价值标签推出因果机制相同。

第三类机制是组合表示。设编码器$E$、组合算子$\circ$和解码器$D$，对训练中未出现的组合$a\star b$，希望$D(E(a)\circ E(b))$满足领域的语法与语义约束。理解“红”和“球”并不必然理解“红球”：系统还需知道颜色修饰对象、球可能被红光照亮而非材质为红、词序和指代怎样绑定。若编码对置换、旋转或图同构具有等变性，这一性质只覆盖声明的变换群；Attention本身也不会自动内化现实世界所有对称性。更严谨的结构保持检验要列出对象映射、组合映射、恒等操作和允许近似误差，再用系统性留出组合、反事实属性交换和关系重绑定测试。所谓函子只是适当抽象，不是从物理信号到语义必然存在且保结构的经验定理。

三类机制可以联合，但不能相互替代。局部结构网络提供对象身份、关系和组合规则，全局分布表示提供相似性、不确定性与连续优化；二者通过任务相关身份映射、绑定接口、耦合消息和读出损失连接。局部吸引可以抵抗噪声，却可能固化偏见；结构映射支持远迁移，却可能因遗漏边界条件产生危险类比；组合接口支持新构造，却可能生成语法合法而事实虚假的结果。因此验收矩阵应同时覆盖域内留出、可控扰动、分布偏移、组合留出、因果干预、置信校准、资源消耗和失败恢复。只有外部任务证据支持时，我们才说系统获得了相应泛化能力。

在HSF--HD中，信息几何因此处于方法层：它为分布状态、正定矩阵、局部投影和参数更新提供明确度量；一般状态动力学还要容纳离散身份、图编辑、目标变化和外部行动。AgentOS可用$h(t)$承载快速分布状态，用$E$保留域迁移证据，用$\Theta$保存经验证的结构映射，SPC产生任务读出，TCE选择更新尺度，CCM监控病态矩阵和分布漂移，Boundary与Offloading控制证据进入和外部保存。这个实例说明一种落地方式，却不把HSF--HD缩减为AgentOS。信息几何真正提供的不是泛化神话，而是一套可说明坐标、度量、投影条件、离散误差与失败边界的计算语言。

%--chp---
